%=============================================================================!
% 2.4  THE THIRD LAW                                                          !
%=============================================================================!
\section{The third law}
\label{sec:ne-third}

Forces do not appear from nowhere: every force on a body is exerted by
some other body. The third law describes what happens at the other end,
and it is what later lets us treat several bodies, and many atoms,
together.

We write $\vb{F}_{AB}$ for the force on body $A$ exerted by body $B$; the
first subscript names the body that feels the force. The third law states
that forces between two bodies come in matched pairs,
\begin{equation*}
   \vb{F}_{AB} = -\vb{F}_{BA} ,
\end{equation*}
equal in size and opposite in direction. Like the second law, this is a
fact of experiment: hook two spring balances together and pull them apart,
and they always show the same reading, however hard or unevenly you pull.
When you push on a wall, the wall pushes back on you with an equal force.
When the Earth pulls a falling ball down, the ball pulls the Earth up with
a force of the same size. A ball of half a kilogram falls, with air
resistance negligible, at the acceleration $g$ of Chapter~1, so by the
second law the Earth's pull on it is $0.5\times9.80665 =
\SI{4.9}{\newton}$, its weight (\cref{sec:ne-gravity}). The Earth's mass
is about \SI{6e24}{\kilogram}, so by the second law again the Earth's
acceleration is
$4.9/(6\times10^{24}) = \SI{8e-25}{\metre\per\second\squared}$, far too
small to notice.

The two forces of a pair always act on different bodies, so they never
cancel in the second law of either body. A book resting on a table shows
the commonest confusion (\cref{fig:ne-third}). The Earth pulls the book
down and the table pushes it up, and these two forces are equal and
opposite, but they are not a third-law pair: both act on the book, and
they are equal only because the book does not accelerate. The partner of
the Earth's pull on the book is the book's pull on the Earth, and the
partner of the table's push on the book is the book's push on the table.

\begin{figure}[htb]
\centering
\begin{tikzpicture}[line cap=round,
   force/.style={-{Stealth[length=5pt]}, line width=1.1pt},
   body/.style={line width=0.6pt, fill=black!8}]
   % the book, the table and the Earth, drawn apart
   \draw[body] (-0.6,2.6) rectangle (0.6,3.1);
   \draw[body] (-2.9,1.2) rectangle (2.9,1.34);
   \draw[line width=0.6pt] (-2.7,1.2) -- (-2.7,0.3) (2.7,1.2) -- (2.7,0.3);
   \fill[black!10] (-3.2,-1.6) arc[start angle=140, end angle=40,
      radius=4.18] -- cycle;
   \node at (0,-1.25) {\small Earth};
   \node at (1.05,2.85) {\small book};
   \node[right] at (2.9,1.27) {\small table};
   % one pair: the table pushes the book up, the book pushes the table down
   \draw[force] (-0.25,3.1) -- (-0.25,3.9)
      node[midway, left] {\small table on book};
   \draw[force] (-0.25,1.34) -- (-0.25,0.4)
      node[midway, left] {\small book on table};
   % the other pair: the Earth pulls the book, the book pulls the Earth
   \draw[force, accent] (0.25,2.6) -- (0.25,1.8)
      node[midway, right] {\small Earth on book};
   \draw[force, accent] (0.25,-0.1) -- (0.25,0.7)
      node[midway, right] {\small book on Earth};
\end{tikzpicture}
\caption{The forces between a book resting on a table, the table and the
Earth, with the three bodies drawn apart for clarity. The black arrows
are one third-law pair, the table's push on the book and the book's push
on the table; the teal arrows are another, the Earth's pull on the book
and the book's pull on the Earth. The two forces on the book itself, one
from each pair, balance because the book does not accelerate. Forces
between the table and the Earth are not drawn.}
\label{fig:ne-third}
\end{figure}

\begin{selfcheck}
The Earth pulls an apple down and the apple pulls the Earth up with an
equal force. Why does the apple fall towards the Earth, rather than the
Earth towards the apple?
\end{selfcheck}
